\documentclass{article}
\usepackage[utf8]{inputenc}

\usepackage{amsmath}
\usepackage{amssymb}
\usepackage{enumitem}
\usepackage{marvosym}
\usepackage{amscd}

\title{Coreflexive Membership}
\author{Adam W. Grant}
\date{January 2021}

\begin{document}

\maketitle
The set comprehension
\begin{equation}
    \{ z : p.z : z\}
\end{equation}
describes a set of scalars.  We can transform this into a set of pairs by adjusting the body of the relation
\begin{equation}
    \{ z : p.z : (z,z)\}
\end{equation}
Now what does it mean for an arbitrary pair $(y,x)$ to be a member of (2)?  Let's calculate\footnote{All hints in the proof name theorems from \cite{10.5555/1963395} except for ``coreflexive notation", which comes from \cite{Oliveira2009}.}:
\begin{itemize}
    \item $(y,x) \in \{ z : p.z : (z,z)\}$
    \item[$\equiv$] $\{\text{Set membership}\}$
    \item[] $\exists(z : p.z : (y,x) = (z,z))$
    \item[$\equiv$] $\{\text{Pair equality}\}$
    \item[] $\exists(z : p.z : y=z \wedge x=z)$
    \item[$\equiv$] $\{\text{Double trading}\}$
    \item[] $\exists(z : x=z : y=z \wedge p.x)$
    \item[$\equiv$] $\{\text{One-point rule}\}$
    \item[] $y=x \wedge p.x$
    \item[$\equiv$] $\{\text{Coreflexive notation}\}$
    \item[] $y \Phi_{p} x$
    \item[\Heart]
\end{itemize}
And so we see that
\begin{equation*}
    (y,x) \in \{ z : p.z : (z,z)\}
\end{equation*}
\begin{equation*}
    y=x \wedge p.x
\end{equation*}
\begin{equation*}
    y \Phi_{p} x
\end{equation*}
have the same meaning.
\bibliographystyle{plain}
\bibliography{refs}
\end{document}
